\section{Discussion and Conclusion}

The two refinements proposed in this paper address distinct failure modes
of Algorithm~\ref{alg:naive} of \cite{hou2021fundamentals}. Iterative
deepening, the more impactful of the two, recovers proofs that the
depth-first baseline misses by getting stuck in deep
quantifier-instantiation chains: it accounts for 49 of the 55 additional
solves observed in our combined prover. The LK$'$ six-class priority
schedule contributes a smaller but consistent additional gain,
concentrated on problems with many atoms where the visible-terms-first
phase pays for itself by avoiding unnecessary fresh-witness generation.

The remaining 70\,\% of unsolved problems are dominated by timeouts
rather than soundness-preserving \emph{Unknown} outcomes, indicating
that the prover would benefit more from sharper search heuristics than
from a larger resource budget. Promising directions include
unification-based instantiation in the manner of
leanTAP~\cite{beckert1995leantap}, term ordering for $\forall L$/$\exists R$
witness selection, and lemma sharing across iterative-deepening
iterations to amortise the re-traversal overhead. Extending the
evaluation to the unprovable problem class would also test the dual
question of how quickly the prover can saturate a non-theorem and report
\emph{NotProvable} rather than time out.

The principal lesson of this study is that small, principled changes to
the rule-application order of a textbook sequent calculus can produce
disproportionate improvements in practical solve rate, even before any
of the heavyweight machinery (resolution, superposition, free-variable
unification) of mature first-order provers is brought to bear.

\subsubsection{Data Availability.}
Source code, benchmark scripts, and the SQLite database backing every
figure in this paper are available at the project repository on GitHub.
